% Unidad R569: introducción, líneas francesas 95--120, pp. 8--9.
% Traducción: OpenAI Codex, GPT-6.1 Sol, esfuerzo Ultra.
Por último, es evidente que un libro sobre geometría algebraica, y sobre todo un libro sobre sus fundamentos, recibe necesariamente la influencia, aunque sea por intermediarios, de matemáticos como O.~Zariski y A.~Weil.
En particular, la \emph{Théorie des fonctions holomorphes} de Zariski~\cite{I-20}, convenientemente flexibilizada gracias a los métodos cohomológicos y completada con un teorema de existencia (cap.~III, \S\S~4 y~5), es, junto con la técnica de descenso expuesta en el cap.~VI, una de las principales herramientas empleadas en este tratado y nos parece una de las más potentes de que se dispone en geometría algebraica.

La técnica general en la que se inscribe puede esbozarse del siguiente modo (se ofrecerá un ejemplo típico en el cap.~IX, al estudiar el grupo fundamental).
Se tiene un morfismo propio (cap.~II) $f:X\to Y$ de una variedad algebraica a otra (más generalmente, de un esquema a otro) que se quiere estudiar en una vecindad de un punto $y\in Y$, con el fin de resolver un problema P relativo a una vecindad de $y$.
Se procede por etapas sucesivas:
\begin{enumerate}
  \item[$1^\circ$] Se puede suponer que $Y$ es afín, de modo que $X$ pasa a ser un esquema definido sobre el anillo afín $A$ de $Y$, e incluso se puede sustituir $A$ por el anillo local de $y$.
  Esta reducción siempre es fácil en la práctica (cap.~V) y nos lleva al caso en que $A$ es un anillo \emph{local}.
  \item[$2^\circ$] Se estudia el problema considerado cuando $A$ es un anillo local \emph{artiniano}.
  Para que siga teniendo efectivamente sentido cuando no se supone que $A$ sea íntegro, a veces es preciso reformular el problema P, y resulta que a menudo se obtiene así una mejor comprensión del problema, de naturaleza «infinitesimal» en esta etapa.
  \item[$3^\circ$] La teoría de esquemas formales (cap.~III, \S\S~3, 4 y~5) permite pasar del caso de un anillo artiniano al de un \emph{anillo local completo}.
  \item[$4^\circ$] Por último, si $A$ es un anillo local cualquiera, la consideración de «secciones multiformes» sobre esquemas adecuados sobre $X$, que aproximan una sección «formal» dada (cap.~IV), permitirá a menudo pasar de un resultado conocido para el esquema

  \oldpage[I]{9}
  deducido de $X$ por extensión de escalares a la completación de $A$, a un resultado análogo para una extensión finita suficientemente sencilla (por ejemplo, no ramificada) de $A$.
\end{enumerate}

Este esbozo muestra la importancia del estudio sistemático de los esquemas definidos sobre un anillo artiniano $A$.
El punto de vista de Serre en su formulación de la teoría local de cuerpos de clases, y algunos trabajos recientes de Greenberg, parecen sugerir que ese estudio podría emprenderse asociando funtorialmente a un esquema $X$ de este tipo un esquema $X'$ sobre el cuerpo residual $k$ de $A$ (supuesto perfecto), de dimensión igual, en los casos favorables, a $n\dim X$, donde $n$ es la longitud de $A$.

En cuanto a la influencia de A.~Weil, baste decir que la necesidad de desarrollar las herramientas para formular con toda la generalidad deseada la definición de la «cohomología de Weil» y para abordar la demostración\footnote{Para evitar cualquier malentendido, precisemos que esta tarea apenas acaba de emprenderse en el momento en que se escribe esta introducción y, por tanto, todavía no ha llevado a la demostración de las conjeturas de Weil.} de todas las propiedades formales necesarias para establecer sus célebres conjeturas en geometría diofántica~\cite{I-19} ha sido una de las principales motivaciones para redactar este tratado, al igual que el deseo de encontrar el marco natural de las nociones y métodos habituales de la geometría algebraica y de dar a los autores la oportunidad de comprender esas nociones y técnicas.

\bigskip
\begin{center}*\quad*\quad*\end{center}
\bigskip
